\documentclass[11pt]{article}
\usepackage[letterpaper,margin=.87in]{geometry}
\usepackage{amsmath,amssymb,booktabs,hyperref,microtype}
\ifdefined\pdfmapfile\pdfmapfile{+cm.map}\pdfmapfile{+cmextra.map}\pdfmapfile{+symbols.map}\fi
\hypersetup{colorlinks=true,urlcolor=blue,pdftitle={The largest active prime at a quadratic modulus cap},pdfauthor={Codex (OpenAI)}}
\setlength{\parindent}{0pt}\setlength{\parskip}{5pt}
\newcommand{\head}[1]{\par\vspace{5pt}{\large\bfseries #1}\par}
\newcommand{\lcm}{\operatorname{lcm}}
\newcommand{\eps}{\varepsilon}
\begin{document}
{\LARGE\bfseries The largest active prime at a quadratic modulus cap}\par
{\large Irreducible covering sets of every exact cardinality}\par
Codex (OpenAI), at the request of Jordan Pollard. Research manuscript for review, 27 November 2026.

A set $A$ of $k$ distinct integers greater than one covers if one residue for each modulus covers all integers. It is irreducible if every proper subset fails for all choices of residues. Write $M=\max A$, $L=\lcm A$ and $P=P^+(L)$. Define $P_\lambda(k)$ as the maximum $P$ among such sets with $M\le\lambda k^2$, or zero for an empty family.

\head{The sharp tradeoff at every size}
For every fixed $\lambda>0$,
\[
 \lim_{k\to\infty}\frac{P_\lambda(k)}k
   =\min\{1,\sqrt{(1-\log 2)\lambda}\}.\tag{1}
\]
Equivalently, for every fixed $0<t<1$, the least $M$ among irreducible $k$-sets with $P\ge tk$ is asymptotic to $t^2 k^2/(1-\log 2)$. At cap $k^2$ the answer in (1) is about $0.55394$. At intermediate caps $k^\sigma$, $1<\sigma<2$, the sharp prime scale is $\sqrt{(1-\log\sigma)k^\sigma}$ (page 8). The transition to $P/k\to1$ occurs at a cap of about $3.25889k^2$.

\head{Exact inverse problem with a prescribed prime}
Let $K(p,T)$ be the fewest moduli in any irreducible set with $p\mid L$ and $M\le pT$. Here $p$ is prime and $T$ is an integer. Define
\[
 b(Q)=\sum_{q^a\parallel Q}a(q-1),\quad
 C(p,T)=\min_{(Q,p)=1,\ \#\{d\mid Q:d\le T\}\ge p} b(Q).
\]
For every sufficiently large prime $p$, simultaneously for every $T\ge p+1$,
\[
                         K(p,T)=p+C(p,T).\tag{2}
\]
There is no example for $T\le p$. The right side is a finite arithmetic optimization over nonincreasing exponent lists. For $p/T\to\alpha\in(0,1)$ let $y_*$ be the least prime for which at least $p$ integers up to $T$ are $y_*$-smooth. Then
\[
                  K(p,T)-p\sim \frac{y_*^2}{2\log y_*}.\tag{3}
\]
The same scale controls how many integers lose coverage when the multiples of $p$ are removed. The greatest possible density of that set of integers, maximizing also over covering assignments, is $\exp(-(1+o(1))y_*)$.

Erd\H os \#1189(ii) asks for the unrestricted minimum $M$ at each $k$. Equation (1) determines its restriction $P\ge tk$; the unrestricted minimum and the sharp unrestricted reciprocal coefficient remain open here. The proofs use the Zn\'am--Simpson period bound, the classical prime fibre observation, the earlier digit construction, and Dickman's smooth-number theorem. The verification files concern finite instances.

\newpage
\head{1. Why the inverse problem is arithmetic}
Every irreducible covering set satisfies $1+b(L)\le k$ (Zn\'am--Simpson). This also applies to the period of a minimal covering subfamily under any given assignment. Fix $p^a\parallel L$. A modulus divisible by $p^a$ has a private point in a covering assignment. Vary only the top base-$p$ digit of that point in CRT coordinates. Every modulus missing that digit misses the entire fibre. At least $p$ different moduli divisible by $p^a$ are required. On division by $p^a$ these give distinct divisors of $Q=L/p^a$ at most $pT/p^a\le T$. This proves $k\ge1+a(p-1)+b(Q)\ge p+C(p,T)$. The first $p$ positive integers coprime to $p$ already reach $p+1$.

For the reverse bound compress the prime support of $Q$ onto the first primes different from $p$. Moving a prime to a smaller unused one decreases cost and maps its divisors injectively to no larger divisors. Arrange exponents in decreasing order. If $u<v$ have exponents $a<b$, swap them. For divisors whose exponent at $v$ exceeds $a$, interchange their exponents at $u,v$; keep other divisors fixed. This is an injection into no larger divisors of the new integer. Every prime power $q^a\parallel Q$ is at most $T$ in a cost minimizer, since an unused top digit can be dropped. We use a compressed minimizer below.

\textbf{Realization criterion.} Suppose the primes of $Q$ are less than $p$, its top powers are at most $T$, and each $q^a\parallel Q$ has $q-1$ distinct divisors of the earlier-prime part of $Q$ at most $pT/q^a$. Suppose at least $p$ divisors of $Q$ are at most $T$. For each $j\le a$ choose the indicated $q-1$ distinct earlier cofactors $d$; include the moduli $dq^j$. Assign the nonzero $q$-digits $b=1,\ldots,q-1$ bijectively, by the congruences
\[
                 x\equiv0\pmod d,\quad x\equiv bq^{j-1}\pmod {q^j}.
\]
Choose $p$ distinct terminal cofactors $e\mid Q$, $e\le T$, including every full prime power of $Q$ (there are fewer than $p$). Include $pe$ with $x\equiv0\pmod e$ and a different residue modulo $p$ for each $e$. All moduli are distinct by their largest prime, its exponent and their cofactor. An integer with a first nonzero digit in the ordered prime digits of $Q$ is covered there; all multiples of $Q$ reach the terminal block. The nonterminal rows miss precisely $Q\mathbb Z$.

The period is $pQ$ and there are $p+b(Q)=1+b(pQ)$ rows. For any proper divisor $D\mid pQ$ at most $b(D)$ rows divide $D$: each surviving digit permits at most $q-1$, and retaining every terminal row forces $pQ\mid D$. A covering subfamily minimized after arbitrary reassignment needs at least $1+b(D)$ rows at its own period $D$. At period $pQ$ it needs every row. This proves irreducibility in the required sense.

\newpage
\head{2. Realizing an optimizer uniformly in the cap}
Here $p$ tends to infinity. If $T\ge2p$, feasibility at $T_0=2p$ gives $C(p,T)=o(p^2/\log p)$ uniformly. Indeed take all prime powers with bases $q\le (2p)^{9/10}$ up to $2p$. Their cost is $O(p^{1.8}\log p)$ and their divisors through $2p$ include all smooth numbers with that cutoff. Dickman's formula $\Psi(x,x^{1/u})/x\to\rho(u)$ gives proportion $1-\log(10/9)>1/2$. The prime number theorem gives $\sum_{q\le x}(q-1)\sim x^2/(2\log x)$. Hence the largest prime in an optimal compressed support is $o(p)$, since all smaller primes except possibly $p$ occur. 

For $p+1\le T<2p$ there is no optimally used prime $v>p$. Such a prime accounts for just one divisor through $T$, namely $v$. Remove it. If fewer than $p$ divisors remain, some element of $\{1,\ldots,p-1,p+1\}$ is missing. Increase an insufficient exponent of one of its prime factors by one. This introduces at least the new prime power itself through $T$ and costs less than $v-1$. This contradicts optimality. We have used that $p$ is odd.

For each prime $q$ in either compressed support the earlier part contains the squarefree product of all smaller primes. For all large $q$ this gives at least $q-1$ divisors at most $2q$. Indeed the squarefree integers through $2q$ have asymptotic count $12q/\pi^2$; deleting those with a prime factor $\ge q$ costs $O(q/\log q)$. For each of the bounded smaller $q$, the earlier product has at least $q-1$ divisors by Bertrand's postulate: if $q$ is the $j$th prime then $q-1\le2^{j-1}$. Their sizes are bounded by an absolute constant for this finite set, hence by $p$ eventually.

The allowance $pT/q^a$ is at least $p$. This handles $q\le p/2$ and thus all primes when $T\ge2p$. Otherwise $p/2<q<p$, $T<2p$, and $a=1$ for large $p$. If $2q^2\le pT$ the allowance still reaches $2q$. In the last case $q>\sqrt{pT/2}>T/2$. Of the at least $p$ divisors through $T$, those with a prime factor $\ge q$ must themselves be primes in $[q,p)$, at most $p-q$ in number. At least $q$ earlier divisors remain, all within the allowance $pT/q\ge T$. This proves the realization criterion in every case and completes (2).

At a fixed small prime the arithmetic lower bound still applies. Checking the criterion on one optimizer suffices to prove equality for that instance. For example at $p=11$ the exact minima for $T=12,16,22,44$ are respectively $28,21,20,18$. Explicit optimal cofactor periods are $2520,240,120,72$. The accompanying program searches exponent lists and constructs the rows and residues.

\newpage
\head{3. Smooth divisors and the size of the excess}
Write $p/T\to\alpha\in(0,1)$. Define $\beta\in(0,1)$ by $\rho(1/\beta)=\alpha$. Dickman's theorem and strict monotonicity of $\rho$ beyond one give $y_*=T^{\beta+o(1)}$. Every feasible cofactor period contains at least $\pi(y_*)$ primes: otherwise compress them onto the first primes and even the unrestricted smooth count is too small. Thus its cost is at least $\sum_{q\le y_*}(q-1)$.

Fix $\eta>0$ and put $z=(1+\eta)y_*$. Use all primes $q\le z$ with exponent one, enhancing those $q\le z^{1/2}$ to their largest power at most $T$. Among the $z$-smooth integers through $T$ this loses at most $T\sum_{n>\sqrt z} n^{-2}=O(T/\sqrt z)$ integers. Including primes $y_*<q\le z$ gains $\gg_{\alpha,\eta}T/\log y_*$ integers: take $q$ times any $y_*$-smooth integer through $T/z$. These sets are disjoint, and Dickman's theorem (or monotone squeezing with fixed exponents) gives a positive smooth proportion at that bound. Thus at least $p$ divisors remain. The extra powers cost $O(z\log T)$ and increase $\log Q$ by $O(\sqrt z\log T)$. The cost is $(1+o(1))z^2/(2\log z)$, and $\log Q\sim z$. Let $\eta\downarrow0$ to get (3). This period satisfies the realization criterion by the squarefree argument on page 3 since $z=o(p)$.

\textbf{The threshold for (1).} The inverse function in (3) crosses $\beta=1/2$ exactly at $\alpha_0=\rho(2)=1-\log2$. If $M\le\lambda k^2$ and along a subsequence $P/k$ exceeds $\sqrt{\alpha_0\lambda}$ by a fixed amount, set $p=P$. Then $p\asymp k$ since $p\le k$, and $p^2/M$ stays above $\alpha_0$ by a fixed amount. Enlarge the cofactor cap if needed to $T=\lceil p/\alpha\rceil$ with fixed $\alpha\in(\alpha_0,1)$ below that ratio. Equations (2)--(3) would force $k-p\ge p^{2\beta-o(1)}$ with $2\beta>1$, an impossibility. This proves the upper bound in (1).

\textbf{Exact cardinality in the opposite direction.} Fix $0<t<\min(1,\sqrt{\alpha_0\lambda})$. Choose a constant $g>1/\alpha_0$ with $gt^2<\lambda$; take a prime $p\sim tk$ and $T=\lceil gp\rceil$. Choose $0<\beta'<1/2$ such that $\rho(1/\beta')>1/g$. Start with the primes up to $y=T^{\beta'}$. Give each prime its first digit, then enlarge its exponent as far as $q^a\le T^{1/4}$ allows. This period $Q_0$ has cost $O(y^2\log T)=o(p)$. It has at least $p$ divisors through $T$: powers first missing at $q\le T^{1/8}$ exceed $T^{1/4}$, costing $O(T^{7/8})$ integers in total. Squares at primes above $T^{1/8}$ cost another $O(T^{7/8})$. The starting smooth count exceeds $p$ by a fixed positive proportion of $T$.

\newpage
\head{Reaching every size and retaining the cap}
Set $r=k-p-b(Q_0)>0$. For every prime $q$ offer one optional digit with cost $q-1$: one more exponent for old primes, the first exponent for new ones. Every integer through the sum of the first $j$ costs is a subset sum. In fact the first cost is one and each successive cost is at most one plus the preceding sum, by induction and Bertrand ($q_{i+1}-1\le2(q_i-1)$). Stop at the first prefix reaching $r$. Its largest prime is $O(\sqrt{r\log(2r)})=o(p)$ by the prime number theorem. Select a subset totaling $r$.

All top powers remain at most $T$. For an old prime an additional power is at most $\max(qT^{1/4},q^2)\le T$; new primes have exponent one. Old primes have all smaller primes available and the allowance is at least $p$, so the squarefree argument applies. For any new prime $q>y$, at least $p$ divisors of $Q_0$ through $T$ are available earlier, hence at least $q-1$ within $pT/q$. The terminal row conditions survive. The realized irreducible set has exactly $k$ moduli, $P=p$ and $M\le pT\le\lambda k^2$ for large $k$. Taking $t$ up to the stated limit proves (1). Inverting the strict inequalities gives the fixed-$t$ minimum statement on page 1.

\head{4. A sharp estimate for loss of coverage}
For any assignment covering all integers remove every row divisible by $p$, and let $H$ be the periodic set now uncovered. With $p/T\to\alpha\in(0,1)$ and $M\le pT$, the maximum density $\Delta(p,T)$ of $H$ over irreducible sets and their assignments satisfies
\[
                       -\log \Delta(p,T)\sim y_* .\tag{4}
\]
The period constructed on page 4 leaves precisely the multiples of $Q$ and $\log Q\sim(1+\eta)y_*$. It gives the required lower bound for the density.

For the upper bound no $p^2$ can divide the period of an irreducible set in this regime: the top-digit bound gives $p^a\le T=O(p)$. At each $x\in H$, varying the $p$ coordinate while keeping the others fixed still misses all the retained rows. At least $p$ distinct cofactors $d\le T$ must therefore hit $x$ in their assigned congruences modulo $d$.

Take $X$ uniform in $H$ over a common period containing all prime coordinates up to $T$ (enlarge by their product if needed). Let $D=\log(1/\operatorname{dens}H)$; the empty case is immediate. For each prime $q$ write $v_q=\max_a\Pr(X=a\bmod q)$. Subadditivity of Shannon entropy on CRT coordinates gives
\[
 D\ge\sum_{q\le T}D_{\rm KL}(X\bmod q\;\|\;\mathrm{Unif}_q).
\]
Indeed $D$ is the deficit against the uniform law on the full period; every fibre of projection to these coordinates has the same size. By grouping the most probable residue versus all others, each marginal deficit is at least $v_q\log(qv_q)-v_q$ (use $(1-v)\log(1-v)\ge-v$).

\newpage
\head{Arithmetic estimates for the entropy bound}
Take $w=y_*^{1/2}$ and $s=(\log T)^{-2}$. For $s\le t\le1$ retain all primes $\le w$ and all larger primes $q\le T$ with $v_q\ge t$. Let $Y(t)$ be the prime at the same rank as the cardinality of this set. Compression and integration of indicators give
\[
 p\le sT+\int_s^1\Psi(T,Y(t))\,dt,\qquad
 \int_s^1Y(t)\,dt\le(1+o(1))(D+O(w)).
\]
For the first inequality, a hit modulo $d$ has probability at most the minimum $v_q$ over its prime factors above $w$, with empty minimum one. At each level compress those allowed prime factors to obtain at most $\Psi(T,Y(t))$ choices of $d$. For the second, $\theta(Y(t))$ is the least log product of that many primes. Integrating this costs at most $\theta(w)+\sum_{q>w,v_q\ge s}v_q\log q$. The entropy bound controls the sum with relative error $o(1)$ since $\log q\gg_\alpha\log T$ and $|\log s|=O(\log\log T)$. The prime number theorem applies uniformly since $Y(t)\ge(1-o(1))w$.

Write $u=1/\beta$, $h=\rho(u-1)>0$, with $\rho(v)=1$ for $0\le v\le1$. For each fixed $a>0$,
\[
 \Psi(T,ay_*)-\Psi(T,y_*)=(h\log a+o(1))\,T/\log y_*.
\]
For $a>1$ count by the largest prime in $(y_*,ay_*]$. For each such $q$ the summand is $\Psi(T/q,q)$, whose ratio to $T/q$ tends uniformly to $h$ by monotonicity and Dickman's theorem. For $a<1$ reverse the interval, with the same proof. Partial summation with the prime number theorem gives the formula, and monotonicity upgrades it to uniformity for $a$ in compact subintervals of $(0,\infty)$. The jump at $y_*$ is at most $T/y_*=o(T/\log T)$.

Consequently for all $Y\in[w/2,T]$,
\[
 \Psi(T,Y)\le p+\frac{T}{\log y_*}\left[h(Y/y_*-1)+o(1)\right].
\]
On compact ratios this uses $\log a\le a-1$. Below $e^{-2}y_*$ use monotonicity at $e^{-2}y_*$. Above a sufficiently large constant multiple, each doubling interval $[v,2v]$ adds at most $T\sum_{v<q\le2v}1/q\ll T/\log y_*$ by Chebyshev. The resulting $O(1+\log a)$ is at most $h(a-1)$ for all sufficiently large $a$. Substitution into the integrated bound yields $\int_s^1Y(t)dt\ge(1-o(1))y_*$, proving (4). This argument also bounds the upper tail for arbitrary choices of one residue per $d\le T$ when at least $p$ must hit.

\head{Sources and remaining question}
The original formulation is at \href{https://www.erdosproblems.com/1189}{Erd\H os Problems 1189}. The period inequality is due to Zn\'am and Simpson; the prime fibre bound is classical. For smooth numbers see A. Granville, \emph{Smooth numbers: computational number theory and beyond} (2008). The companion every-size manuscript gives the earlier digit and terminal-cofactor construction. Equation (1) leaves the regime $P=o(k)$ open for the unrestricted minimum. In particular the estimate does not show the unrestricted minimum has order $k^2$; previous constructions achieve $O(k\log k)$. The independent role of medium-sized primes at such smaller caps remains to be determined.
\newpage
\head{5. Without a cap: the exact prime spectrum}
There is also an exact formula for the largest possible prime factor of a modulus at prescribed cardinality, if $M$ is free. Set
\[
 F(v)=\min\left\{\sum_{i=1}^s a_i(q_i-1):
 \prod_{i=1}^s(1+a_i)\ge v,\ a_1\ge\cdots\ge a_s\ge1\right\},
\]
where $q_i$ are the consecutive primes from 2. For every sufficiently large prime $p$, the sizes of irreducible sets having largest prime factor exactly $p$ are precisely the integers $k\ge p+F(p)$. Hence the largest possible $P$ at every sufficiently large $k$ is the largest prime $p$ with $p+F(p)\le k$.

Indeed the lower bound follows from the top-digit count, the period bound and compression with the cap removed. Taking the first $\lceil\log_2 p\rceil$ primes to exponent one costs $O((\log p)^2\log\log p)$. Every prime in an optimizer is therefore below $p$ eventually. Its divisor count reaches $p$. With no bound on the sizes of the earlier divisors, the criterion on page 2 holds by the squarefree product and Bertrand argument. Increase the exponent of 2 by any nonnegative integer to add exactly that many rows; available cofactors and their count only increase. The period continues to have largest prime $p$. For all sufficiently large $k$ the displayed maximum exceeds the finite exceptional range by this very construction.

The excess has a sharp leading coefficient. With $g_j=\log(1+1/j)$ and $A=\sum_{j\ge1}g_j^2$,
\[
                  F(v)\sim\frac{(\log v)^2\log\log v}{2A}.
\]
The $j$th exponent increment at prime $q$ costs $q-1$ and increases the log divisor count by $g_j$. Select all increments $q-1\le xg_j$. This is a feasible prefix in each exponent and across primes. Write $V(x)$ for its total gain and $B(x)$ for cost. The prime number theorem at each fixed $j$ gives
\[
 V(x)\sim A x/\log x,\qquad B(x)\sim A x^2/(2\log x).
\]
For $j\le\sqrt x$ the summands divided by their respective scales are $O(g_j^2)$ by Chebyshev estimates (the cutoffs are $\gg\sqrt x$). For $\sqrt x<j=O(x)$, bound counts and sums of primes by the counts and sums of integers up to $1+x/j$. These tails are respectively $O(\sqrt x+\log x)$ in gain and $O(x^{3/2})$ in cost. No larger $j$ qualifies. For a fixed budget $B(x)$ no selection can give greater gain than $V(x)$: all selected increments have gain/cost at least $1/x$ and all excluded ones have smaller ratio. Taking $x=(1\pm\eps)\log v\log\log v/A$ brackets the required cost; let $\eps$ decrease to zero. This statement concerns the necessary excess at a specified prime and does not require information about prime gaps below a given cardinality.

\head{Verification and credits}
The repository includes an exact constructor, an independent small-cap optimization using all cofactor LCMs, and Lean 4.33.1 files. At $p=11,T=16$ Lean proves that the exhibited 21 distinct moduli cover precisely when all are retained, with arbitrary reassignment. It also checks the finite exponent table excluding cofactor cost below 10. The general period lower bound and the analytic theorems in this manuscript are written proofs. AI authorship is indicated on page 1. The literature credits above apply throughout; priority and external review are pending.

\newpage
\head{6. Between the near-linear and quadratic caps}
The argument gives a continuous range of sharp unrestricted extremal results for the period prime, including cases where it grows sublinearly in cardinality. For $1<\sigma<2$,
\[
 \max_{|A|=k,\;M\le k^\sigma} P
      \sim \sqrt{(1-\log\sigma) k^\sigma}.
\]
All maxima here run over genuine irreducible sets. At the smaller cap $M\le\lambda k\log k$ with fixed $\lambda>1$, the corresponding maximum is asymptotic to $\sqrt{\lambda k\log k}$.

For the power-cap upper bound set $X=k^\sigma$. If $p=P$ exceeds $\sqrt{\rho(\sigma)X}$ by a fixed factor, then $p\asymp\sqrt X$, since always $p^2\le M$. Choose a fixed $\alpha>\rho(\sigma)$ strictly below $p^2/X$ on this subsequence, and enlarge the cofactor cap to $\lceil p/\alpha\rceil$. In (3) the resulting $\beta$ is greater than $1/\sigma$, so $k\ge K(p,T)\ge p^{2\beta-o(1)}$ is impossible. Here $\rho(\sigma)=1-\log\sigma$.

For the lower bound fix $0<t<\sqrt{\rho(\sigma)}$, take a prime $p\sim t\sqrt X$ and $T=\lfloor X/p\rfloor$. Choose $\beta'<1/\sigma$ so that $\rho(1/\beta')>t^2$. Use precisely the base period $Q_0$ on page 4, with cutoff $y=T^{\beta'}$. Its cost is $o(k)$ and at least $p$ divisors survive through $T$. The old primes are $o(p)$. For padding to exact size use optional digits at old primes $q\le\sqrt T$ and all new primes $q>y$. This omits at most the one interval $(\sqrt T,y]$. If the interval is present, the available preceding costs sum to $(1+o(1))T/\log T$, which exceeds $2y$. Hence the first cost beyond the gap is at most one plus this sum. Subsequent costs obey the same completeness induction as on page 5. If the interval is absent there is no gap. All selected top powers are at most $T$. The omitted costs total $O(y^2)=o(k)$, so the largest prime of the first eligible prefix reaching $r=k-p-b(Q_0)\sim k$ is $(1+o(1))\sqrt{k\log k}$, which is $o(p)$. Every new prime has $p$ earlier cofactors through $T$; every old prime still meets the squarefree criterion. The construction has $k$ rows and maximum at most $pT$. Let $t$ tend to $\sqrt{\rho(\sigma)}$.

For $X=\lambda k\log k$ fix $1/\sqrt\lambda<t<1$ and choose $\beta'<1$ with $\rho(1/\beta')>t^2$. Exactly the same bounds hold with $y^2=o(k)$. The largest padding prime is now $(1+o(1))\sqrt{k\log k}<p$. This last strict inequality suffices in place of $o(p)$ for new primes; the old primes remain $o(p)$. The proof realizes every sufficiently large $k$. Taking $t$ up to one attains the elementary upper bound. These estimates are compatible with earlier unrestricted constructions having $M\le(1+o(1))k\log k$. The question of whether the unrestricted minimum can be significantly smaller requires information about simultaneous essentiality beyond a single prescribed prime.

\newpage
\head{7. The reciprocal maximum at the first possible size}
For a prescribed large prime $p$, consider all irreducible sets with $p\mid L$ at the first possible size $k=p+F(p)$. Write $S_{\rm first}(p)$ for their greatest reciprocal sum. Define $W(y)=\sum_{q\le y} H_{q-1}/(q-1)$ over primes, and $Y_p=(\log 2/A)\log p\log\log p$, with $A$ as on page 7. Then
\[
 S_{\rm first}(p)=W(Y_p)+o(1)
 =\log\log p+(1+\gamma)\log\log\log p+C_0+\log(\log 2/A)+o(1),
\]
where $C_0=\lim_{z\to\infty}(W(z)-\log z-\gamma\log\log z)$ is the classical prime-sum constant from the companion manuscript. This is a sharp reciprocal maximum in a prescribed-prime subfamily of genuine unrestricted-residue irreducible sets. The unrestricted maximum at the same cardinality can grow on the larger scale $\log k$.

\textbf{The optimizer contains all small divisors.} The lower bound on page 2 forces $L=pQ$ with $b(Q)=F(p)$, $\tau(Q)\ge p$, and zero defect. Every optimal $Q$ has consecutive prime support from 2, since compression over a skipped smaller prime strictly reduces cost. Let $l$ be its largest prime; all its primes are smaller than $p$ as on page 7. For each $q<l$ with exponent $a_q$, remove one digit of $l$. This reduces $\tau$ by at most a factor of two. Adding $a_q+1$ digits of $q$ doubles the old factor $a_q+1$, so optimum cost forces $(q-1)(a_q+1)\ge l-1$. Since $q^{a_q+1}\ge q(a_q+1)>l-1$, all integers through $l-1$ divide $Q$.

\textbf{Location of the last prime.} Put $x=(\log p)\log\log p/A$. The cost of an optimum is $(1+o(1))B(x)$ and its gain $G=\log\tau(Q)$ is at least $\log p$. Relative to the threshold set of digits on page 7, its loss in the linear functional gain minus cost divided by $x$ is
\[
 V(x)-G+(F(p)-B(x))/x
 =\sum_{(q,j)\text{ in exactly one set}} |g_j-(q-1)/x|
 =o(\log p).
\]
Here nonnegativity and the stated upper bound give the final equality. If $l/(x\log2)$ differed from one by a fixed amount, an interval of primes of length comparable with $x$ would have its first digits on the wrong side, each a fixed distance from the threshold. The prime number theorem gives loss $\gg x/\log x\asymp\log p$. Thus $l/Y_p\to1$ uniformly in every minimizer.

\textbf{Both bounds for the sum.} In the digit construction, at each stage $q^j$ with $1\le j\le a_q$, use all $d=1,\ldots,q-1$. These divide the earlier part of $Q$. Their sum over the stages is $\sum_{q\mid Q}H_{q-1}(1-q^{-a_q})/(q-1)=W(l)-o(1)$. Indeed for $q\le\sqrt l$ the first missing power is at least $l$ and the total omitted tail is $O(\log l/\sqrt l)$; above $\sqrt l$ use $\sum_{n>\sqrt l}(1+\log n)/n^2=o(1)$. For every competing cover the zero-defect digit matching bound of the companion manuscript gives at most this same unextended finite prime-digit sum, plus $H_{p-1}/p+1/M=o(1)$. Finally $W(l)-W(Y_p)=o(1)$ by the prime number theorem and $l/Y_p\to1$. The evaluation of $W$ follows by $H_{q-1}=\log q+\gamma+O(1/q)$ and the classical Mertens prime sums.

\newpage
\head{8. Many indispensable prime sections in one cover}
The cofactor period can serve multiple prescribed primes simultaneously. Suppose $Q_0$ is a base as in the proof of (1), with at least $P$ divisors through $T=\lfloor X/P\rfloor$ and all its prime bases below the prescribed primes. Append distinct prescribed primes $p_1<\dots<p_s=P$ with $P\le T$ to exponent one. At each nonterminal prescribed prime choose $p_i-1$ earlier divisors from the base. Reserve $P$ for the terminal packet, including all mandatory prime powers of the resulting earlier period. They are at most $T$ and fewer than $P$ in number, since their bases are different primes less than $P$. Thus the terminal set can include them and still have $P$ distinct cofactors through $T$. Primes introduced for cardinality padding may be processed in increasing order together with the prescribed primes. Each new prime above the base cutoff has the needed earlier divisors in $Q_0$. The proof on page 2 gives arbitrary-reassignment irreducibility of the entire set. Its size is one plus the sum of all digit costs.

\textbf{Complete finite-dimensional profile at quadratic caps.} For $M\le\lambda k^2$, sort the distinct period primes in decreasing order, using zeros after exhaustion. For every fixed $s$, the possible limiting lists of the first $s$ primes divided by $k$, attainable along every sufficiently large exact size, are exactly
\[
 t_1\ge\cdots\ge t_s\ge0,\qquad
 \sum_i t_i\le1,\qquad t_1^2\le(1-\log2)\lambda .
\]
Necessity follows from page 1 and the additive period budget. For sufficiency first take strict budget and cap inequalities, and at least one positive $t_i$. Choose distinct primes $p_i\sim t_i k$ for the positive entries, also in case of ties by the prime number theorem. Use the base, cutoff, and cofactor allowance of the proof of (1) at the largest such prime (one can take $T=\lceil gP\rceil$ there). The total cost already assigned is $k\sum_i t_i+o(k)$. Fill the positive remainder by the same complete optional costs from primes $o(k)$, disjoint from all prescribed primes. No padding prime is macroscopic. Boundary points, including the zero vector, follow by a diagonal sequence of attainable strict interior approximants, choosing the approximation index to tend to infinity slowly with $k$. In particular the largest possible $s$th prime divided by $k$ tends to $\min(1/s,\sqrt{\lambda(1-\log2)})$.

\textbf{Number of large primes at an intermediate power cap.} For $X=k^\sigma$, $1<\sigma<2$, and $0<t<\sqrt{1-\log\sigma}$, the greatest possible number of distinct period primes $\ge t\sqrt X$ is $(1+o(1))k/(t\sqrt X)$. Above that threshold there are eventually none (no assertion here at equality). The upper count is the period budget. For the lower, prescribe $\lfloor(1-\eps)k/(t\sqrt X)\rfloor$ primes in $[t\sqrt X,(t+\delta)\sqrt X]$ for fixed sufficiently small $\delta>0$. There are enough since $X/(k\log X)\to\infty$. Their sum is below $(1-\eps/2)k$ if $\delta$ is small. Use page 8's base with $T=\lfloor X/P\rfloor$, choosing a fixed exponent $\beta'<1/\sigma$ with $\rho(1/\beta')>(t+\delta)^2$. The base costs $o(k)$; exact-size padding uses primes $O(\sqrt{k\log k}+y)=o(\sqrt X)$ and avoids all prescribed primes. Uniform smooth abundance follows by monotonicity at the worst endpoint. Let $\eps\downarrow0$.

\newpage
\head{The sharp count at caps near the known constructions}
Fix $\lambda>1$, let $X=\lambda k\log k$, and fix $0<t<1$. Define $N_t(A)$ to be the number of distinct primes dividing the period that exceed $t\sqrt X$. Then
\[
 \max_{|A|=k,\,M\le X}N_t(A)
 \sim 2\left(\min\{1,\sqrt{t^2+1/\lambda}\}-t\right)
                      \frac{\sqrt X}{\log X}.
\]
Unlike the power-cap case, a narrow interval at the bottom of the admissible range can have too few primes. Both availability and total cost enter the constant.

\textbf{Necessity.} For fixed $0<a<b\le1$, the number of primes in $(a\sqrt X,b\sqrt X]$ is $(2(b-a)+o(1))\sqrt X/\log X$, and their total digit cost is $(b^2-a^2+o(1))X/\log X$. All active primes are at most $\sqrt X$. For a fixed count, the lowest-cost choice uses the smallest primes above $t\sqrt X$. The period inequality bounds their total cost by $k=(1/\lambda+o(1))X/\log X$. The two prime estimates give the claimed upper bound.

\textbf{Sufficiency.} Write $z_*=\min(1,\sqrt{t^2+1/\lambda})$. Choose $\max(t,1/\sqrt\lambda)<z<z_*$, then prescribe \emph{all} primes in $(t\sqrt X,z\sqrt X]$. Their largest $P\sim z\sqrt X$ satisfies $T=\lfloor X/P\rfloor\ge P$ for large $k$. Choose $\beta'<1$ with $\rho(1/\beta')>z^2$. The previous truncated-power base at $y=T^{\beta'}=o(\sqrt X)$ costs $o(k)$ and gives at least $P$ available cofactors. The assigned cost from the prescribed primes is $(z^2-t^2+o(1))X/\log X<k$. To fill the remaining cardinality use the optional slots below $t\sqrt X$, omitting the old primes in $(\sqrt T,y]$ when necessary. Their total available cost is $(t^2+o(1))X/\log X$ since $y^2=o(k)$. This strictly exceeds the needed remainder because $z^2>1/\lambda$. The subset-sum argument across the omitted interval is the one on page 8. Thus exact size $k$ is attained. Every added prime is below $P$, and the shared-period construction on page 10 gives $M\le X$. Now let $z\uparrow z_*$.

For instance when $\lambda=2$ and $t=1/2$, the asymptotic maximum is $(\sqrt3-1)\sqrt X/\log X$. These examples show that many large primes can be simultaneously indispensable at the scales of interest to the minimum-modulus problem. A bound requiring their small cofactor periods to be disjoint would fail.

\end{document}
